% =====================================================================
%  数学建模竞赛论文 — 主文件
%  编译方式:make pdf (latexmk -xelatex)
%  章节内容按节拆分在 sections/ 下,队员同时写不同文件避免冲突
%  图片统一放 figures/,由 code/src/*_fig.py 生成
% =====================================================================
\documentclass[12pt]{article}

\usepackage{geometry}
\geometry{a4paper, top=2.5cm, bottom=2.5cm, left=2.8cm, right=2.8cm}

% ---------- 中文环境 ----------
\usepackage{iftex}
\ifXeTeX
  \usepackage{fontspec}
  % 优先使用国赛要求的字体(标题黑体、正文宋体)。
  % 检测 Windows 系统字体目录里是否有 SimSun:
  %   有 -> 手动指定中文字体;没有(如 Linux 队友机器) -> 回退 Fandol 开源字体。
  % 用 \IfFileExists 而不是 \ifFontExistsTF,兼容老版本 fontspec(MiKTeX 旧版)。
  \IfFileExists{C:/Windows/Fonts/simsun.ttc}
    {\usepackage[fontset=none]{ctex}
     \setCJKmainfont{SimSun}[AutoFakeBold=true]
     \setCJKsansfont{SimHei}}
    {\usepackage[fontset=fandol]{ctex}} % Fandol 是 TeX Live/MiKTeX 自带的开源中文字体
\else
  \usepackage[UTF8]{ctex}
\fi

% ---------- 图表 ----------
\usepackage{graphicx}
\usepackage{booktabs}    % 三线表
\usepackage{caption}     % 图表标题
\captionsetup{font=small, labelfont=bf}
\usepackage{float}
\usepackage{subcaption}

% ---------- 数学 ----------
\usepackage{amsmath, amssymb, amsthm}
\usepackage{bm}
\numberwithin{equation}{section}

% ---------- 代码 ----------
\usepackage{listings}
\usepackage{xcolor}
\lstset{
  basicstyle=\small\ttfamily,
  breaklines=true,
  frame=single,
  backgroundcolor=\color{gray!8},
  columns=flexible
}

% ---------- 超链接 ----------
\usepackage[hidelinks]{hyperref}
% 参考文献采用 thebibliography 手动编号方案(国赛最稳妥,不依赖外部工具),
% 条目写在 sections/08-references.tex 中。

% ---------- 论文信息 ----------
\title{题目名称\\ \large 摘要:一句话概括工作}
\author{
  队员一 \quad 队员二 \quad 队员三\\
  \small(学校名称 学院名称,省份 城市 邮编)
}
\date{\today}

\begin{document}

\maketitle

% ========== 摘要 ==========
\begin{abstract}
  针对问题一,建立了\ldots模型,采用\ldots算法求解,得到\ldots结论。
  针对问题二,\ldots。
  针对问题三,\ldots。

  \textbf{关键词:}关键词一;关键词二;关键词三;关键词四
\end{abstract}

\tableofcontents

% ========== 正文各节 ==========
% =====================================================================
%  01 问题重述
%  在比赛开始后,把题目原文摘录/概括到这里。
%  通常分两小节:1.1 问题背景,1.2 问题提出。
% =====================================================================
\section{问题重述}

\subsection{问题背景}

(概括题目给出的背景材料,用自己的话重述,不要照抄原文。)

\subsection{问题提出}

\begin{itemize}
  \item 问题一:\ldots
  \item 问题二:\ldots
  \item 问题三:\ldots
\end{itemize}

% =====================================================================
%  02 模型假设
%  写出建模时做的合理假设,每条假设在建模中会被用到,
%  尽量每条都能对应到 04 模型建立里的某个式子/符号。
%
%  符号说明表:由队长(Claude)统一维护,全队以此为准。
%  建模手写公式、代码手写代码、论文手排版都用同一套符号。
%  新符号出现时,更新此表并推送,大家在群里同步。
% =====================================================================
\section{模型假设}

\begin{enumerate}
  \item 假设一:……
  \item 假设二:……
  \item 假设三:……
\end{enumerate}

\newpage
\section*{符号说明}

\begin{table}[htbp]
\centering
\caption{主要符号说明}
\begin{tabular}{lll}
\toprule
符号 & 含义 & 单位 \\
\midrule
$\alpha$   & 针肋宽度比(针肋直径/通道宽度) & — \\
$\beta$    & 歧管深高比(歧管层/微通道层深高比) & — \\
$n$        & 单个歧管单元内沿流向的针肋排数 & — \\
$R_{th}$   & 总热阻 & K/W \\
$\Delta P$ & 压降 & Pa \\
$\sigma_T$ & 温度非均匀性指标 & K \\
$T_{max}$  & 芯片最高温度 & K \\
$T_{in}$   & 冷却液入口温度 & K \\
$Q$        & 热流密度/总热功率 & W/m$^2$ 或 W \\
$h$        & 对流换热系数 & W/(m$^2$$\cdot$K) \\
$A$        & 换热面积 & m$^2$ \\
$D_h$      & 通道水力直径 & m \\
$u$        & 流速 & m/s \\
$\rho$     & 冷却液密度 & kg/m$^3$ \\
$c_p$      & 冷却液比热容 & J/(kg$\cdot$K) \\
$\mu$      & 冷却液动力黏度 & Pa$\cdot$s \\
$f$        & 沿程阻力系数 & — \\
$Re$       & 雷诺数 & — \\
\bottomrule
\end{tabular}
\end{table}

% 说明:以上为预填常用符号,比赛过程中随建模进度增删。
% 新增符号时:告知队长 → 队长更新本表 → 推送,全队同步。

% =====================================================================
%  03 问题分析
%  每个问题用一段连续文字概述解决思路(不分步骤),
%  只讲"怎么做",不写具体公式、不重述问题内容。
%  分步骤的详细建模见第 4 节,以加粗 Step N 引出。
% =====================================================================
\section{问题分析}

\subsection{问题一的分析}

本文基于传热学与流体力学理论\cite{White2016,Yang2006}建立流固耦合控制
方程组,定义热阻、压降与温度非均匀性三项性能指标;进而提取雷诺数、
普朗特数与结构参数比等主导无量纲量,确认附件数据处于层流工况,为后续
代理模型输入选取提供机理依据;在此基础上,依据达西摩擦、局部阻力叠加
与热阻串联机理构造压降幂律关联式与热阻二次结构关联式,由附件数据最小
二乘标定;最后固定其余参数,定量统计各结构参数对三项指标的影响方向与
幅度,验证机理规律并论证三项指标综合评价的合理性。

\subsection{问题二的分析}

对 84 组样本按 $8:2$ 留出法划分训练/测试集,输入特征做 Z-score 标准化;
考虑到样本量小且映射高度非线性,选用高斯过程回归\cite{Rasmussen2006},
以 Constant$\times$RBF+White 复合核建立多输入--单输出(MISO)映射,
用极大似然估计优化超参数;最后以决定系数、均方根误差与平均绝对百分比
误差评估拟合与泛化能力,确认代理模型可替代 CFD 求解器。

\subsection{问题三的分析}

以 GPR 代理模型为目标函数,在样本支撑域内建立三指标同时最小化的多目标
优化模型(热阻-压降负相关 $r=-0.684$ 表明解集为 Pareto 前沿);采用
NSGA-II\cite{Deb2002} 全局搜索 Pareto 前沿,L-BFGS-B\cite{Byrd1995,Chen2005}
局部精修,并对排数整数约束离散修正;最终以 Min--Max 无量纲化消除量纲
差异,用均等权重理想点法\cite{Hwang1981} 将多目标转化为贴近理想解的
单目标决策,唯一确定综合最优设计方案并回代验证。

\subsection{问题四的分析}

在权重单纯形上采样全场景偏好,建立由偏好权重到最优结构参数的敏感性
映射,量化权重波动对方案的驱动作用;以各场景定制最优为基准定义后悔值,
采用最小最大后悔值准则\cite{Savage1951} 求取最坏偏好下综合性能损失最小
的鲁棒方案;最后以权重空间变异系数与最大后悔值对比鲁棒方案与等权综合
方案,论证鲁棒设计的工程价值。

\subsection{问题五的分析}

对加工误差(结构参数)与工况波动(流量、入口温度、发热功率)两类扰动进行
概率分布表征,建立 $\pm5\%$ 综合摄动模型;结构参数层由代理模型计算局部
归一化敏感性系数与 Sobol 全局敏感性指数\cite{Sobol2001},工况参数层基于
无量纲标度关系一阶近似;采用拉丁超立方采样与蒙特卡洛传播\cite{McKay1979}
生成性能响应分布,以变异系数、95\% 置信上限与越界失效概率判定稳定性;
最后将综合最优方案与鲁棒方案一并代入扰动传播,对比评估两类方案的工程
稳定性。

% =====================================================================
%  04 模型建立
%  论文的核心章节。每个问题一个小节,写明:
%   1. 建模思路(为什么用这个模型)
%   2. 符号定义(引用第 2 节的符号表)
%   3. 公式推导
%   4. 模型的解/算法(对应 code/src/ 里的代码)
%
%  公式编号自动带章节前缀,如 (4.1)。
% =====================================================================
\section{模型建立}

\subsection{问题一的模型}

(建模思路:……)

设 $x$ 为……,目标函数为
\begin{equation}
  \min f(x) = \sum_{i=1}^{n} w_i x_i,
  \label{eq:obj}
\end{equation}
约束条件为
\begin{equation}
  \sum_{i=1}^{n} x_i \le C, \quad x_i \ge 0, \ i=1,\ldots,n.
  \label{eq:cons}
\end{equation}

求解方法:采用……算法(详见 \texttt{code/src/} 中对应脚本),流程图如图~\ref{fig:flow} 所示。

\subsection{问题二的模型}

(建模思路:……)

\subsection{问题三的模型}

(建模思路:……)

% =====================================================================
%  05 模型求解与结果分析
%  本文件不再由 main.tex 引入。
%  各问题的模型建立与求解内容已合并到 sections/04-models.tex 中。
% =====================================================================

% =====================================================================
%  06 模型的评价及推广(检验、灵敏度分析、优缺点)
%  结构:6.1 模型检验 → 6.2 灵敏度分析(指向问题五) → 6.3 优缺点
%
%  占位符约定:\textbf{(待回填:……)} 是给编程手留的数值空位,
%  求解脚本跑完后逐个替换,不要删除。
% =====================================================================
\section{模型的评价及推广}

\subsection{模型检验}

\subsubsection{代理模型精度检验}

问题二建立的 GPR 代理模型采用留出法(训练集与测试集按 $8:2$ 划分,
独立测试样本不参与训练)评估拟合与预测精度,三项性能指标的检验结果
见表~\ref{tab:model-check}。由表可见,三个模型的决定系数 $R^2$ 均
接近于 1,均方根误差与平均绝对百分比误差均处于极低量级,表明代理模型
对仿真数据的拟合与预测精度高,可代替 CFD 求解器用于后续的多目标优化
与不确定性传播分析。

% 待回填:以下数字由 code/src/solve_q2.py 运行输出
\begin{table}[h]
\centering
\caption{GPR 代理模型精度检验结果(测试集)}
\label{tab:model-check}
\begin{tabular}{lccc}
\toprule
性能指标 & $R^2$ & $RMSE$ & $MAPE$ \\
\midrule
热阻 $R_{th}$      & \textbf{(待回填)} & \textbf{(待回填)} & \textbf{(待回填)} \\
压降 $\Delta P$    & \textbf{(待回填)} & \textbf{(待回填)} & \textbf{(待回填)} \\
温度非均匀性 $\Delta T$ & \textbf{(待回填)} & \textbf{(待回填)} & \textbf{(待回填)} \\
\bottomrule
\end{tabular}
\end{table}

\subsubsection{机理规律与仿真数据的一致性检验}

问题一推导的参数影响规律与附件 2 样本数据的控制变量分析结果对照:
针肋宽度比对热阻的影响在样本范围内呈"先降后升"的非单调变化,各组
控制变量下极小值均出现在 $\alpha \approx 0.2$ 附近;歧管深高比增大
使流量分配更均匀,温度非均匀性总体下降且影响幅度较小;针肋排数增大
使压降总体升高,且增幅随宽度比增大而加大,均与机理推导一致,表明机理
模型对系统行为的刻画符合仿真规律。定量吻合程度(平均相对偏差等)待
求解结果回填。

\subsubsection{优化结果的可行性检验}

将问题三所得综合最优方案回代 GPR 代理模型验证目标函数值,与优化
过程中预测值的相对偏差小于 \textbf{(待回填:偏差百分比)} $\%$,
验证了优化结果的可行性与自洽性。

\subsection{灵敏度分析}

结构参数与工况参数对性能指标的敏感性已在问题五中系统完成:结构参数
扰动对三项指标的影响由训练好的代理模型严格刻画,采用 Sobol 全局敏感
性指数(一阶主效应、总效应)与局部归一化敏感性系数定量评估,并基于
蒙特卡洛不确定性传播给出变异系数、$95\%$ 置信上限与越界失效概率;
工况参数(流量、入口温度、发热功率)的影响基于标称设计点处的梯度
线性化近似评估,详见 5.5 节。分析表明,系统对 \textbf{(待回填:最敏感
参数与结论)} 的扰动最为敏感,该结论为加工公差控制与运行工况约束提供
了定量依据。

\subsection{模型的优缺点}

\textbf{优点:}
\begin{itemize}
  \item 机理与数据驱动融合:问题一建立 CFD 机理模型,问题二以 GPR
        代理模型逼近高维非线性映射,兼顾物理可解释性与拟合精度;
  \item GPR 除给出点预测外还提供预测方差,为后续优化与不确定性
        传播提供了天然的概率框架;
  \item 采用 NSGA-II 全局搜索与 L-BFGS-B 局部精化结合的混合优化
        策略,并对针肋排数做整数离散修正,兼顾解的全局性与收敛精度;
  \item 决策层完整:理想点法消除量纲差异给出唯一综合方案,Minimax
        后悔值准则降低方案对偏好权重的依赖,变异系数进一步量化了
        方案的稳定性;
  \item 敏感性分析体系完备:局部(LSA)与全局(Sobol)结合,并给出
        越界失效概率等工程可操作指标,直接服务于制造公差与运行
        工况约束设计。
\end{itemize}

\textbf{缺点:}
\begin{itemize}
  \item GPR 代理模型的预测精度依赖训练样本的网格覆盖范围,外推至
        样本空间之外或极端工况时精度难以保证;
  \item 机理模型作了均匀体热源、忽略热辐射等简化假设,与实际芯片
        封装结构的传热过程存在一定偏差;
  \item 蒙特卡洛不确定性传播的收敛性依赖采样规模,参数维度升高时
        计算代价增大;
  \item 问题五中工况参数(流量、入口温度、发热功率)对性能的影响基于
        标称设计点处的梯度线性外推近似,仅在 $\pm 5\%$ 小扰动范围内
        有效,扰动幅度增大时误差随之增大;
  \item Minimax 后悔值准则偏保守,在极端偏好场景下综合性能可能
        略逊于针对性优化方案。
\end{itemize}

% =====================================================================
%  07 模型推广
%  简述模型在其他场景的应用潜力,一两段即可,不要凑字数。
% =====================================================================
\section{模型推广}

(模型可推广到\ldots场景。例如:本文建立的\ldots方法同样适用于\ldots)

% 注:国赛正文(不含摘要页)限 20 页,各部分务必精炼,图约 12-15 张为宜。

% =====================================================================
%  08 参考文献
%  用 thebibliography 手动编号(国赛最稳妥的方案,不依赖外部工具)。
%  格式:
%   [N] 作者. 题名[J]. 期刊名, 年, 卷(期): 页码.     (期刊)
%   [N] 作者. 题名[M]. 出版地: 出版社, 年.            (书籍)
%  正文中用 \cite{key} 引用;编号按正文首次出现顺序排列。
% =====================================================================
% 注意:cumcmthesis 重定义的 thebibliography 环境自带 \section*{\refname} 标题,
% 此处不再重复写 \section{参考文献} 避免双标题。
\begin{thebibliography}{99}

\bibitem{Tuckerman1981}
Tuckerman~D~B, Pease~R~F~W. High-performance heat sinking for VLSI[J].
IEEE Electron Device Letters, 1981, 2(5): 126--129.

\bibitem{White2016}
White~F~M. Fluid Mechanics (8th ed.)[M]. New York: McGraw-Hill, 2016.

\bibitem{Yang2006}
杨世铭, 陶文铨. 传热学(第4版)[M]. 北京: 高等教育出版社, 2006.

\bibitem{Rasmussen2006}
Rasmussen~C~E, Williams~C~K~I. Gaussian Processes for Machine Learning[M].
Cambridge: MIT Press, 2006.

\bibitem{Hwang1981}
Hwang~C~L, Yoon~K. Multiple Attribute Decision Making: Methods and
Applications[M]. Berlin: Springer-Verlag, 1981.

\bibitem{Deb2002}
Deb~K, Pratap~A, Agarwal~S, et al. A fast and elitist multiobjective genetic
algorithm: NSGA-II[J]. IEEE Transactions on Evolutionary Computation, 2002,
6(2): 182--197.

\bibitem{Byrd1995}
Byrd~R~H, Lu~P, Nocedal~J, et al. A limited memory algorithm for bound
constrained optimization[J]. SIAM Journal on Scientific Computing, 1995,
16(5): 1190--1208.

\bibitem{Chen2005}
陈宝林. 最优化理论与算法(第2版)[M]. 北京: 清华大学出版社, 2005.

\bibitem{Savage1951}
Savage~L~J. The theory of statistical decision[J]. Journal of the American
Statistical Association, 1951, 46(253): 55--67.

\bibitem{Sobol2001}
Sobol~I~M. Global sensitivity indices for nonlinear mathematical models and
their Monte Carlo estimates[J]. Mathematics and Computers in Simulation,
2001, 55(1--3): 271--280.

\bibitem{McKay1979}
McKay~M~D, Beckman~R~J, Conover~W~J. A comparison of three methods for
selecting values of input variables in the analysis of output from a computer
code[J]. Technometrics, 1979, 21(2): 239--245.

\bibitem{Tao2001}
陶文铨. 数值传热学(第2版)[M]. 西安: 西安交通大学出版社, 2001.

\end{thebibliography}


\end{document}
